\uuid{vh59}
\titre{Partie entière}
\auteur{bodin ; hanani ; mzari}
\organisation{exo7}
\licence{CC BY-NC-SA 4.0}
\datecreate{2020-01-20}
\niveau{L1}
\module{Analyse}
\chapitre{Propriétés de R}
\sousChapitre{Propriétés des nombres réels}
\difficulte{1}
\theme{Propriétés de nombres réels, partie entière}
\exo7qcmid{lille-1.reels.009}

\contenu{
\begin{question}
\qclassification{120.03}
\qtags{motcle=partie entière}
Notation : $E(x)$ désigne la partie entière du réel $x$.
Quelles sont les assertions vraies ?
\begin{answers}
    \bad{$E(7,9) = 8$}

    \good{$E(-3,33) = -4$}

    \bad{$E(\frac{5}{3}) = 5$}

    \bad{$E(x)=0 \implies x=0$}
\end{answers}
\begin{explanations}
La partie entière de $x$ est le plus grand entier, inférieur ou égal à $x$.
\end{explanations}
\end{question}
}
